% NOTE: revisit this abstract once the pending experiments (Section~\ref{sec:results-discussion})
% are finalized -- it currently reflects the committed methodology and the results available
% at drafting time.

On-board processing promises to overcome the downlink latency and bandwidth constraints that limit the usefulness of Earth observation (EO) data for time-critical applications such as flood and wildfire response. Geospatial Foundation Models (GeoFMs) are attractive candidates for this role because their large-scale, multimodal pre-training yields semantic feature extractors that generalize across tasks with few labeled examples. However, GeoFMs are pre-trained on archived data from missions such as Sentinel-2, and their behavior when exposed to a new sensor is unknown. This thesis investigates whether TerraMind, a Vision Transformer (ViT)-based GeoFM, produces domain-invariant latent representations when applied to real imagery from the European Space Agency's $\Phi$-sat-2 CubeSat mission, and whether that invariance can be induced through domain adaptation and preserved through compression to a hardware-deployable architecture.

Using a co-registered triplet dataset of real $\Phi$-sat-2, real Sentinel-2, and physics-simulated $\Phi$-sat-2 imagery, we show that the domain gap between the two sensors is measurable and concentrated in early encoder layers, driven primarily by differing spectral response functions rather than by spatial resolution mismatch. We further show that naively fine-tuning TerraMind on physics-degraded data does not reduce this gap relative to an undegraded baseline, motivating an explicit, paired-contrastive domain-invariance objective that exploits the co-registration directly. This mechanism is applied during knowledge distillation of TerraMind into a lightweight convolutional neural network (CNN) suitable for the $\Phi$-sat-2 on-board processor, addressing the incompatibility of ViT architectures with space-qualified hardware. The resulting pipeline is evaluated on real-data land-cover classification and on a physics-degraded flood-segmentation benchmark, and compared against $\Phi$satNet, a bespoke on-orbit GeoFM developed for the same mission, to assess whether adapting an existing large-scale pre-trained model offers a label-efficiency advantage over training a mission-specific model from scratch. The results position deliberate, encoder-level domain-invariance training -- rather than incidental exposure to degraded data -- as a necessary condition for realizing the generalization promise of GeoFMs in resource-constrained, on-board Earth observation settings.
